%\documentclass[15pt]{article}
%\input{1.all_header_standalone}
%\begin{document}
We adapt the forest algorithm from \cite{cho_multi-stage_2023} to single stage treatment settings. We provide a high-level algorithm of the forest in Algorithm S1. Briefly, we use the Kaplan-Meier estimator to estimate the marginal survival probability of a node (or leaf) $\Node$, where each individual’s failure time information is given by $\{(X_i, \delta_{i})\}^{n_{\Node}}_{i=1}$ where \(n_{\Node}\) is the size of node \(\Node\), $X_i$ is the observed failure time, and $\delta_i$ is the event indicator for any cause: 
\begin{equation}
\label{eqn:KMest}
\survfunchat(t\mid \Node) = \prod_{s \geq 0} \left( 1 + \frac{\sum_{i \in \Node} dN_{i}(s)}{\sum_{i \in \Node} Y_{i}(s)} \right),
\end{equation}
where $N_i(t) = I(X_i \le t, \delta_i = 1)$ and $Y_i(t) = I(X_i \ge t)$.
The Kaplan-Meier estimator is used in both tree partitioning and terminal node survival probability estimation. 

When the optimization criterion is the truncated mean survival time (\( \phi = \phi^\mu \)), the forest algorithm splits a node into two candidate daughter nodes, \( \Node_1 \) and \( \Node_2 \) by using the mean difference test to maximize the difference in truncated mean survival times:
\begin{equation}
\label{eqn:MDtest}
MD(\Node_1, \Node_2) = \left| \int_{t > 0} (t \wedge \tau) d\widehat{S}(t \mid \Node_1) - \int_{t > 0} (t \wedge \tau) d\widehat{S}(t \mid \Node_2) \right|.
\end{equation}
This splitting rule identifies the best split as \( (\Node_1, \Node_2) = \arg \max_{\Node_1,\Node_2} MD(\Node_1, \Node_2) \). 

The recursive partitioning process continues as long as the size of each node remains above the threshold \( n_{\min} \). The survival probability estimate of the terminal node \( \Node \) of the \( w \)-th (\( w = 1, 2, \ldots, n_{\text{tree}} \)) survival tree is given by the Kaplan-Meier estimator \( \widehat{S}_{w}(t \mid \Node) \). Letting $h = (a,Z)$, the random survival forest estimate is then given by
\[
\widehat{S}(t \mid h) = \frac{1}{n_{\text{tree}}} \sum_{w=1}^{n_{\text{tree}}} \sum_{l=1}^{\#\{\Node_{\{w\},l} : l\}} 1(h \in \Node_{\{w\},l}) \widehat{S}(t \mid \Node_{\{w\},l}),
\]
where \( \#\{\Node_{\{w\},l} : l\} \) is the number of terminal nodes of the \( w \)-th tree.

\begin{algorithm}[H]
  \caption{The random survival forest estimator}
  \label{alg:RSF}
  \begin{flushleft}
    \begin{algorithmic}
  \State \textbf{Input:} $\{(X_i, \delta_i, Z_i, A_i)\}_{i=1}^n$
  \State \textbf{Output:} $\widehat{S}(\cdot|Z,A=a)$
  \State Result: random survival forest estimate $\widehat{S}(\cdot|Z,A)$
  \State Parameters: the number of trees $n_{\text{tree}}$, the minimum terminal node size $n_{\text{min}}$
  \State For \underline{treatment arm} $a \in \mathcal{A}$ \textbf{do} 
  \begin{itemize}[label={}]
  \item For \underline{tree $w = 1, 2, \ldots, n_{\text{tree}}$} \textbf{do}
  \begin{itemize}[label={}]
  \item Partition the feature space with the input data $\{(I(t\leq X_i), \delta_{ji}, Z_i)\}_{i:a_i=a}$ via MD test \cref{eqn:MDtest} until node size $< 2n_{\text{min}}$
  \item Obtain the terminal node survival probability estimate $\widehat{S}{\{w\}}(t | a, Z)$ via Kaplan-Meier estimator \cref{eqn:KMest}
  \end{itemize}
  
  \item Obtain the forest estimator $\widehat{S}(\cdot | Z, A = a) = \frac{1}{n_{\text{tree}}} \sum_{w=1}^{n_\text{tree}} \widehat{S}_{\{w\}}(t |a, Z)$
  \end{itemize}
\end{algorithmic}
\end{flushleft}
\end{algorithm}

%\end{document}